\chapter{State of the Art and Design Choices}

\section{Summary}
\label{sec:state-of-art-summary}

This chapter documents the state of the art that emerged from our research. It covers the technology choices, architecture decisions, and synchronization strategies adopted for this project, with citations and critical assessment for each. Where Chapter~2 surveyed the landscape and identified the gap, this chapter explains how the research translated into concrete design decisions. Table~\ref{tab:research-mapping} shows how each research finding maps to a component in the system.

\begin{table}[htbp]
\centering
\caption{Research-to-Implementation Mapping}
\label{tab:research-mapping}
\begin{tabular}{p{4.5cm} p{4cm} p{5cm}}
\toprule
\textbf{Research Finding} & \textbf{Source} & \textbf{System Component} \\
\midrule
Zero-config discovery & RFC 6762/6763 & Go mDNS broadcaster \\
\midrule
Causal event ordering & Lamport (1978) & Per-file Lamport counters \\
\midrule
Partition tolerance & CAP theorem (2002) & Offline-first availability \\
\midrule
Conflict resolution & Bayou (1995) & LWW + local backups \\
\midrule
Optimistic replication & Saito \& Shapiro (2005) & State-transfer architecture \\
\midrule
Delta sync (future) & rsync (1996) & Planned block-level sync \\
\midrule
File change detection & inotify / FSEvents & C++ kernel watchers \\
\midrule
Data integrity & SHA-256 / TCP & C++ hashing + Go TCP pools \\
\bottomrule
\end{tabular}
\end{table}

\section{Zero-Configuration Peer Discovery}

Multicast DNS (mDNS) and DNS-Based Service Discovery (DNS-SD) \cite{rfc6762, rfc6763} are the protocols that allow devices on a local network to find each other without a central server. They are the reason a printer shows up on your laptop's Wi-Fi list without anyone typing in its IP address.

Our Go discovery component broadcasts a service record under the \texttt{\_p2psync.\_tcp} domain and continuously browses the subnet for other instances. When a new peer appears, its address is resolved automatically. Users never see an IP address, a device ID, or an invitation code.

Why mDNS and not something else? During the research phase, Nizar looked into how Syncthing handles discovery. Syncthing uses a combination of local discovery and global relay servers. The local part is effective, but the global relay adds complexity and an external dependency. Since our target environment is a campus Wi-Fi where all team members are on the same subnet, mDNS gives us what we need without any external infrastructure \cite{syncthing}.

There is a known limitation. Some enterprise Wi-Fi networks block multicast traffic. On EPITA's campus network, multicast works without issues. For networks that block it, a fallback would be needed, but adding a relay server would contradict the serverless design goal. For now, we accept this limitation and document it.

\section{Event Ordering with Lamport Clocks}

Every file in the system carries a Lamport counter in its metadata \cite{lamport1978}. When a peer modifies a file, it increments the counter and includes the new value in the synchronization message. A receiving peer compares the incoming counter against its local value. If the incoming number is higher, the file is newer.

This is the same mechanism described in Section~2.2.1, applied at the per-file level. The counter is stored in the local SQLite metadata database alongside the file's SHA-256 hash and path.

One design question came up during implementation: should the counter be global (one counter per peer, shared across all files) or per-file? A global counter would increment even when unrelated files change, leading to unnecessary updates. A per-file counter only increments when that specific file changes, which keeps the update traffic proportional to actual changes. After discussing it in a team meeting, the per-file approach won.

The limitation is unchanged: Lamport clocks cannot detect concurrent edits. If Arihant edits \texttt{report.pdf} on his laptop while Huizhe edits the same file on his, and neither has seen the other's change, the timestamps do not tell the system that a conflict exists. The LWW tiebreaker picks one version, and the backup mechanism saves the other. Fidge's \citeyear{fidge1988} vector clocks would detect the concurrency, but at the cost of carrying an array of counters that must be updated whenever the peer group changes.

\section{Availability Over Consistency}

The CAP theorem \cite{brewer2002} forced a clear decision: prioritise availability. Peers must be able to edit files freely whether or not other peers are reachable. Reconciliation happens when connectivity returns.

In practice, this means the system is ``eventually consistent'' in the sense described by Vogels \citeyear{vogels2009}. A peer that goes offline for an hour can accumulate many local changes. When it reconnects, the Go daemon exchanges metadata with other peers and pulls in any files that have newer timestamps. If a conflict exists, LWW resolves it automatically and the backup directory preserves the losing version.

The risk of eventual consistency is that two peers may briefly see different versions of the same file. For our use case, this is acceptable. A project team member seeing a file from ten seconds ago is fine. A system that prevents editing during a Wi-Fi dropout is not.

\section{Last-Write-Wins with Version Backups}

Conflict resolution uses LWW, following the approach from Bayou \cite{bayou1995}. Before the C++ daemon writes an incoming file to disk, it copies the existing local version into a \texttt{.versions/} subdirectory, tagged with a timestamp. This ensures that LWW never causes permanent data loss.

Saito and Shapiro \citeyear{saito2005} classify our approach as ``state transfer with rollback support.'' State transfer means sending entire file states (as opposed to operation transfer, which sends edit deltas). It is simpler to implement and debug, though it uses more bandwidth.

The team considered three alternatives before settling on LWW:

\begin{enumerate}
    \item \textbf{Manual conflict resolution} (like Git): rejected because it requires user intervention, which defeats the goal of invisible synchronization.
    \item \textbf{CRDTs} \cite{shapiro2011}: rejected for the initial version due to implementation complexity, but the metadata schema is designed to support register-based CRDTs in a future iteration.
    \item \textbf{Rename-on-conflict} (like Dropbox's ``conflicted copy''): considered, but it creates confusing duplicate files. LWW with backups achieves the same safety without cluttering the working directory.
\end{enumerate}

\section{File Change Detection}

Polling a directory for changes is wasteful. It consumes CPU cycles even when nothing has changed, and the polling interval adds latency. Modern operating systems provide kernel-level event APIs: inotify on Linux \cite{inotify7} and FSEvents on macOS \cite{fsevents}.

The C++ daemon registers watchers on the synchronization directory and all its subdirectories. When a file is created, modified, or deleted, the kernel pushes an event. The daemon debounces rapid sequences of events (common when saving a file triggers multiple write operations) and forwards a consolidated update to the Go daemon through IPC.

This gives near-instant detection. In informal tests on Arihant's MacBook, the delay between saving a file and the Go daemon receiving the metadata update was consistently under 50 milliseconds.

\section{Data Integrity and Network Transport}

Every file has a SHA-256 hash \cite{fips1804} computed by the C++ daemon using memory-mapped I/O. The hash serves two purposes: detecting changes (if the hash is different, the file has changed) and verifying integrity after transfer (the receiver recomputes the hash and compares).

Network transport uses TCP \cite{rfc9293}. The Go daemon maintains a pool of TCP connections with discovered peers. When a file transfer is needed, the Go daemon establishes the connection and hands the socket to the C++ daemon, which streams the file data directly to disk. TCP's built-in flow control and error recovery mean we do not need to implement our own retransmission logic.

Whole-file transfer is the initial design. The rsync algorithm \cite{rsync1996} could reduce bandwidth by transferring only changed blocks, but implementing it correctly adds significant complexity. The architecture has a clean separation between the networking layer (Go) and the storage layer (C++), so adding block-level transfers later would only require changes in the C++ daemon.

